\section{Trabalhos Relacionados}

A literatura que antecede esta revisão organiza-se em duas vertentes, e o posicionamento deste estudo decorre da lacuna entre elas.

\noindent\textbf{Revisões de agentes e de Engenharia de Software com IA.} A primeira vertente reúne revisões e visões que mapeiam o uso de modelos de linguagem e de agentes na Engenharia de Software. Hou et al. conduzem uma revisão sistemática ampla do uso de modelos de linguagem em tarefas de Engenharia de Software, organizada por etapa do ciclo de vida e por técnica \cite{hou2024llmse}; He, Treude e Lo revisam a arquitetura interna e a cooperação de sistemas multiagente baseados em modelos de linguagem, delineando um roteiro para a consolidação desses sistemas \cite{JundaHe2025}; e Hassan et al. propõem a visão de uma Engenharia de Software nativa de IA, com um roteiro de desafios para o que chamam de Engenharia de Software 3.0 \cite{AhmedEHassan2024}. Essas contribuições são valiosas para entender a anatomia dos agentes e a direção do campo. Mas operam em um nível de abstração que classifica modelos, componentes cognitivos e técnicas por tarefa, ou que projeta cenários futuros. Nenhuma caracteriza os frameworks operacionais que, no presente, organizam o trabalho de desenvolvimento sobre um agente em nível de processo e de repositório.

\noindent\textbf{Propostas orientadas a processo.} A segunda vertente é formada pelos próprios frameworks que estruturam o processo. O desenvolvimento orientado por especificação articula a especificação como contrato vivo na era dos assistentes de codificação \cite{DeepakBabuPiskala2026}; a distinção entre codificar por intuição e codificar de forma agêntica formaliza um espectro de autonomia com implicações práticas \cite{RanjanSapkota2025}; e frameworks como MetaGPT, ChatDev e Reversa instanciam, cada um a seu modo, a organização do ciclo por papéis, artefatos e validação \cite{SiruiHong2023,ChenQian2023,Macedo2026}. Cada uma dessas propostas, porém, avalia a si mesma de forma isolada, sem um referencial comum que permita compará-las.

\noindent\textbf{Posicionamento.} Entre uma literatura de revisão que descreve agentes de forma abstrata e uma literatura de propostas que avalia cada framework isoladamente, falta uma revisão sistemática que tome os frameworks operacionais como objeto, com uma lente de Engenharia de Software, e que confronte a estrutura desses frameworks com a evidência disponível por meio de um referencial explícito. É essa a contribuição deste artigo: uma taxonomia de seis dimensões confrontada com a evidência de um corpus de 41 estudos, uma síntese temática que expõe os eixos de tensão do campo e uma agenda de pesquisa orientada a processo. O diferencial é mensurável pela unidade de análise: enquanto Hou et al. organizam o uso de modelos de linguagem por atividade do ciclo de vida e por técnica \cite{hou2024llmse}, e He et al. organizam os sistemas multiagente por componente arquitetural e de cooperação \cite{JundaHe2025}, esta revisão toma o framework operacional como unidade e classifica os 41 estudos por seis dimensões de processo, com seis temas indutivos que cobrem 100\% do corpus. Onde aquelas revisões respondem o que modelos e agentes fazem por tarefa, esta responde como o processo de desenvolvimento é estruturado sobre o agente. Diferente das comparações de produto voltadas a praticantes, o foco aqui é o corpus de literatura e a estrutura que dele emerge, não a recomendação de uma ferramenta específica.
